\ifx\allfiles\undefined
\documentclass[12pt, a4paper, oneside, UTF8]{ctexbook}
\def\path{../config}
\input{\path/package.tex}

% 定理环境设置
\input{\path/theorem.tex}

\input{\path/custom.tex}

\input{\path/cover_package} 

% 封面页设置
\def\myTitle{量子力学：理论与应用
}
\def\mySubtitle{Quantum Mechanics: Theory and Application}
\def\myAuthor{Max}
\def\myCreateDateCover{2026年10月11日}
\def\myUpdateDateCover{ \today}
\def\myDateForeword{}
% 前言页设置
\def\myForeword{前言}
\def\myForewordText{

}
\def\mySubheading{}
\begin{document}
% \input{\path/cover_text.tex}

%%%%%%%% 前言页 %%%%%%%%

\newpage
\thispagestyle{empty}
\begin{center}
    \Huge\textbf{\myForeword}
\end{center}
\myForewordText
\begin{flushright}
    \begin{tabular}{c}
        \myDateForeword
    \end{tabular}
\end{flushright}

% 目录页
\newpage
\pagestyle{plain}
\setcounter{page}{1}
\pagenumbering{Roman}
\tableofcontents

% 正文页
\newpage
\pagenumbering{arabic}
\setcounter{chapter}{-1}
\setcounter{page}{1}

\pagestyle{fancy}
% \fancyhead[L]{量子力学笔记}
% \fancyhead[C]{中间页眉}
% \fancyhead[R]{右侧页眉}
% \fancyfoot[C]{\thepage}
\fancyhf{}
% 设置页眉
\fancyhead[L]{\leftmark}    % 左侧：文章标题（\leftmark 通常是章标题）
\fancyhead[C]{ }                       % 中间：留空（或自定义内容）
\fancyhead[R]{\rightmark}   % 右侧：节标题（\rightmark 通常是节标题）


% 设置页脚（示例：中间显示页码）
\fancyfoot[C]{\thepage}

\renewcommand{\headrulewidth}{0.4pt}
\renewcommand{\footrulewidth}{0pt}








\else
\fi

\chapter{态和力学量的表象}
\begin{examct}{}{}
    态的表象、力学量算符的矩阵表示、矩阵力学、幺正变换、Dirac符号、线性谐振子及占有数表象
\end{examct}
\begin{examrqm}{}{}
    \begin{enumerate}
        \item 掌握表象变换、幺正变换、Dirac符号，会求力学量算符的本征值和本征态，掌握线性谐振子的代数解法
        \item 能综合运用矩阵力学的基本知识解决问题
    \end{enumerate}
\end{examrqm}

\section{态的表象}

\subsection{各表象能量本征方程}
\begin{enumerate}
    \item 坐标表象：
    \begin{equation}
        \left[-\frac{\hbar^2}{2m}\frac{\mathrm{d}^2}{\mathrm{d}x^2}+V(x)\right]\psi(x)=E\psi(x)
    \end{equation}

    \item 动量表象：
    \begin{equation}
        \begin{matrix}
            & \displaystyle -\frac{p^2}{2m}\varphi(p)+\int V_{pp^{\prime}}\varphi(p^{\prime})\mathrm{d}p^{\prime}=E\varphi(p) \\[10pt]
            & \displaystyle V_{pp^{\prime}}=\frac{1}{2\pi\hbar}\int V(x)e^{\mathrm{i}(p^{\prime}-p)x/\hbar}\mathrm{d}x
        \end{matrix}\label{eq:EnergyEigenequationinMomentumRepresentation}
    \end{equation}
\end{enumerate}








\section{力学量算符的矩阵表示}
\section{矩阵力学}
\section{幺正变换}
\section{Dirac符号}
\section{线性谐振子及占有数表象}

 


\newpage
\section{附录}
\subsection{动量表象能量本征方程推导}
推导式\eqref{eq:EnergyEigenequationinMomentumRepresentation}.由不进入表象的态矢量的能量本征方程出发
\begin{equation}
    \hat{H}\ket{\psi}=(\hat{T}+V)\ket{\psi}=E\ket{\psi}
\end{equation}
等式两边分别插入恒等算符
\begin{equation}
    \int\ket{p^{\prime}}\bra{p^{\prime}}\mathrm{d}p^{\prime} = 1 
\end{equation}
得到
\begin{equation}
    \int(\hat{T}+V)\ket{p^{\prime}}\langle p^{\prime}|\psi \rangle\mathrm{d}p^{\prime}=E\int\ket{p^{\prime}}\langle p^{\prime}|\psi \rangle\mathrm{d}p^{\prime}
\end{equation}
又$\varphi(p^{\prime})=\langle p^{\prime}|\psi \rangle$ 得
\begin{equation}
    \int(\hat{T}+V)\ket{p^{\prime}}\varphi(p^{\prime})\mathrm{d}p^{\prime}=E\int\ket{p^{\prime}}\varphi(p^{\prime})\mathrm{d}p^{\prime}
\end{equation}
两边同时左乘 $\bra{p}$ 得
\begin{equation}
    \int \bra{p}(\hat{T}+V)\ket{p^{\prime}}\varphi(p^{\prime})\mathrm{d}p^{\prime}=E\int \langle p|p^{\prime} \rangle \varphi(p^{\prime})\mathrm{d}p^{\prime}
\end{equation}
\begin{equation}
    \int \bra{p}\hat{T}\ket{p^{\prime}}\varphi(p^{\prime})\mathrm{d}p^{\prime} +\int \bra{p}V\ket{p^{\prime}}\varphi(p^{\prime})\mathrm{d}p^{\prime}=E\int \langle p|p^{\prime} \rangle \varphi(p^{\prime})\mathrm{d}p^{\prime}
\end{equation}
有克罗内克符号$\delta_{pp^{\prime}}=\langle p|p^{\prime} \rangle$ ,$\hat{T}$的本征方程 $\hat{T}\ket{p^{\prime}}=\frac{p^2}{2m}\ket{p^{\prime}}$

\begin{equation}
    \int \frac{p^2}{2m} \delta_{pp^{\prime}} \varphi(p^{\prime}) \mathrm{d}p^{\prime} + \int V_{pp^{\prime}} \varphi(p^{\prime})\mathrm{d}p^{\prime} = E \int \delta_{pp^{\prime}} \varphi(p^{\prime})\mathrm{d}p^{\prime}
\end{equation}
由$\delta$函数的挑选性 $\int \delta_{pp^{\prime}} \varphi(p^{\prime})\mathrm{d}p^{\prime} = \varphi(p)$ 得

\begin{equation}
    \boxed{
    \frac{p^2}{2m}\varphi(p) + \int V_{pp^{\prime}} \varphi(p^{\prime})\mathrm{d}p^{\prime} = E \varphi(p)
    }
\end{equation} 
该式即为动量表象的能量本征方程.

其中 $V_{pp^{\prime}}$ 
\begin{equation}
        V_{pp^{\prime}} = \bra{p}V\ket{p^{\prime}}
\end{equation}
插入恒等算符
\begin{equation}
    \int\ket{x}\bra{x}\mathrm{d}x = 1 
\end{equation}
得到
\begin{equation}
    \begin{aligned}
        V_{pp^{\prime}}  =& \bra{p}V\ket{p^{\prime}} \\[5pt]
        =& \iint \langle p|x \rangle \bra{x} V \ket{x^{\prime}} \langle x^{\prime}|p^{\prime}\rangle \mathrm{d}x\mathrm{d}x^{\prime}\\[5pt]
        =& \iint \frac{1}{\sqrt{2\pi\hbar}}e^{-\mathrm{i}p^{\prime}x/\hbar} \ V(x^{\prime})\langle x|x^{\prime}\rangle \ \frac{1}{\sqrt{2\pi\hbar}}e^{\mathrm{i}px/\hbar}\ \mathrm{d}x\mathrm{d}x^{\prime}\\[5pt]
        =& \frac{1}{2\pi\hbar} \iint e^{-\mathrm{i}p^{\prime}x/\hbar} V(x^{\prime}) \delta_{xx^{\prime}}\ e^{\mathrm{i}px/\hbar} \ \mathrm{d}x\mathrm{d}x^{\prime}\\[5pt]
        =& \frac{1}{2\pi\hbar} \int e^{-\mathrm{i}p^{\prime}x/\hbar} V(x) \ e^{\mathrm{i}px/\hbar} \ \mathrm{d}x \\[5pt]
        =& \frac{1}{2\pi\hbar} \int V(x) e^{-\mathrm{i}(p^{\prime}-p)x/\hbar}\ \mathrm{d}x \\
    \end{aligned}
\end{equation}
所以
\begin{equation}
    \boxed{
        V_{pp^{\prime}}= \frac{1}{2\pi\hbar} \int V(x) e^{-\mathrm{i}(p^{\prime}-p)x/\hbar}\ \mathrm{d}x
    }
\end{equation}

\ifx\allfiles\undefined
\end{document}
\fi